%Predložak za seminarski rad%
%Zaključak%
%
%
\phantomsection
\addcontentsline{toc}{chapter}{Zaključak}
\chapter*{Zaključak}
\label{chap:conc}
%
U ovom radu realizovan je softverski, odnosno
\foreignlanguage{english}{bit-banged}, \foreignlanguage{english}{UART}
protokol na mikrokontroleru \foreignlanguage{english}{ATmega328P}, bez
upotrebe njegovog ugrađenog hardverskog \foreignlanguage{english}{UART}
modula. Na osnovu opkoda i tačnog broja ciklusa korištenih instrukcija
izvedena je opšta formula za broj ponavljanja petlje za kašnjenje, koja
je zatim primijenjena na punu bit periodu i na poravnanje na sredinu
bita prilikom prijema, čime je teorijski određen broj ponavljanja
petlje potreban za ostvarivanje zadate brzine prijenosa.

Ispravnost implementacije provjerena je funkcionalnim testiranjem na
fizičkoj razvojnoj ploči Arduino Uno R3, uz korištenje alata \emph{avra}
za asembliranje i alata \emph{avrdude} za programiranje mikrokontrolera.
Komunikacijom uspostavljenom preko terminalskog alata \emph{tio}
utvrđeno je da program ispravno prima i vraća nazad svaki poslani
karakter, čime je potvrđena ispravnost implementiranih potprograma za
slanje i prijem podataka u punom krugu.

Testiranje je, međutim, izvedeno isključivo na ovaj funkcionalan način,
bez upotrebe osciloskopa ili logičkog analizatora, tako da stvarno
trajanje pojedinačnog bita na predajnoj liniji nije nezavisno izmjereno.
Kao pravac daljeg rada preporučuje se upravo takva empirijska provjera,
kao i proširenje implementacije mogućnošću podešavanja brzine prijenosa
i detekcijom greške okvira, čime bi se ostvarena rješenja mogla porediti
sa ponašanjem hardverskog \foreignlanguage{english}{UART} modula pod
istim uslovima.
\newpage
